\documentclass[11pt,a4paper]{article}
\usepackage[T1]{fontenc}
\usepackage[utf8]{inputenc}
\usepackage[english]{babel}
\usepackage{amsmath,amssymb,booktabs}
\usepackage[margin=26mm]{geometry}
\usepackage[colorlinks=true,urlcolor=blue,linkcolor=blue]{hyperref}
\title{Supplement S2\\Computational Diagnostics}
\author{Andrea Comparini\\\small A curious amateur with proof obligations}
\date{9 October 2026}
\begin{document}
\maketitle

\noindent Accompanying \emph{Covariogram Uniqueness with One Regular Body in
$\mathbb R^3$: A Candidate Proof via Area Measures and Fourier--Laplace Zeros}.
The computations described here are diagnostics. The mathematical argument
in the paper does not use numerical agreement as a premise.

\section{A scalar test of the spectral identities}
The file \path{ricerca/verifica_media_zeri.py} uses only the Python standard
library and the test profile
\[
 A(t)=(1-t^2)(1+t/5),\qquad -1\le t\le1,
 \qquad F(\zeta)=\int_{-1}^1A(t)e^{i\zeta t}\,dt .
\]
Its endpoint slopes are $a=8/5$ and $b=12/5$, its width is $2$, and the
model zero height is
$\gamma=\tfrac12\log(b/a)=0.202732554054\ldots$.
This scalar test probes signs, constants and finite differences; it is
not a uniqueness test on a class of three-dimensional bodies.

The script computes $10\,000$ approximations to nonzero zeros by Newton
iteration, starting from $(2m+1)\pi/2+i\gamma$, $m=1,\ldots,10\,000$.
It applies Newton's method to an explicit numerator $-\zeta^2F(\zeta)$
and uses the corresponding analytic derivative. These computations do
not certify that the resulting points are zeros, or that every zero has
been found.

The values on the imaginary axis are evaluated using the polynomial
endpoint expansions and the formula
\[
 \int_0^2s^k e^{-\tau s}\,ds
 =\frac{k!}{\tau^{k+1}}
   \left(1-e^{-2\tau}\sum_{j=0}^k\frac{(2\tau)^j}{j!}\right).
\]
The formula is exact mathematically; its implementation uses floating-point
arithmetic. As the endpoints are $-1$ and $1$, the centered quantity
$B(\tau)=D(\tau)+(\alpha+\beta)\tau$ coincides with
$D(\tau)=\log F(i\tau)-\log F(-i\tau)$.

\section{The finite-difference quantity used in the paper}
The main paper uses $H(\tau)=D(2\tau)-2D(\tau)$.
For an approximate zero $\zeta_j=x_j+iy_j$, the truncated spectral
calculation uses
\[
 \sum_{j=1}^{10\,000}
 \bigl(R_{x_j,y_j}(2\tau)-2R_{x_j,y_j}(\tau)\bigr),\qquad
 R_{x,y}(\tau)=\log\frac{x^2+(\tau-y)^2}{x^2+(\tau+y)^2}.
\]
The following output was reproduced while preparing this supplement.
The last column is the signed difference between the truncated spectral
sum and the direct Laplace calculation.

\begin{center}
\begin{tabular}{rrrr}
\toprule
$\tau$ & Direct $H(\tau)$ & Truncated sum & Difference\\
\midrule
20  & 0.345166366871 & 0.345166366738 & $-1.332\cdot10^{-10}$\\
100 & 0.393050736275 & 0.393050719631 & $-1.664\cdot10^{-8}$\\
500 & 0.402968515298 & 0.402966436253 & $-2.079\cdot10^{-6}$\\
\bottomrule
\end{tabular}
\end{center}
The theoretical endpoint limit is
$\log(b/a)=0.405465108108\ldots$.
The increasing discrepancy in the last column is reported without
correction or suppression. The calculation omits the infinite tail,
and neither that tail nor rounding errors are enclosed by certified
interval bounds.

\section{Other diagnostics retained from the notebook}
The same script compares the centered logarithm $B(\tau)$ with the
truncated sum of $R_{x_j,y_j}(\tau)$. It also adds the model-tail estimate
\[
 -\frac{4\gamma}{\pi}
   \arctan\frac{\tau}{\pi(10\,000+1)} .
\]
This correction approximates a tail integral; it is not an error
certificate for the true zero sum. The observed values are:
\begin{center}
\begin{tabular}{rrrr}
\toprule
$\tau$ & Direct $B(\tau)$ & Truncated sum & Error after model tail\\
\midrule
20  & $-0.365053362686$ & $-0.364889050321$ & $-5.262\cdot10^{-14}$\\
100 & $-0.397180683717$ & $-0.396359124554$ & $-2.641\cdot10^{-13}$\\
500 & $-0.403800388239$ & $-0.399692925271$ & $-1.327\cdot10^{-12}$\\
\bottomrule
\end{tabular}
\end{center}
The script additionally evaluates $J(\tau)=B(\tau)-\tau B'(\tau)$,
which belonged to an earlier route in the notebook.
The finite-difference proof in the present paper does not use that derivative.

The separate file \path{ricerca/verifica_scambi.py} checks polynomial
identities, cutoff derivatives and rational inequalities for explicit
examples investigated earlier. Its assertions use exact rational
arithmetic; some displayed illustrative values use floating-point
trigonometric functions. All assertions passed in the current run.
For example, the rational convexity margin reported by the script is
$35781657829/40000000000>0$.
These examples are retained in the research notebook and are not premises
of the candidate theorem in the paper.

\section{Repetition and source records}
From the repository root, with Python 3 available:
\begin{verbatim}
python3 ricerca/verifica_media_zeri.py
python3 ricerca/verifica_scambi.py
\end{verbatim}
Neither script requires third-party Python packages. The source files and
the run outputs are provided separately; the latter are saved under
\path{paper/diagnostics/}. Last-digit floating-point results can depend on
the runtime and platform.

These diagnostics should be distinguished from the kernel-checked
conditional theorems in Supplement S1. They supply inspectable checks of
specific calculations, not an independent validation of the geometric or
spectral transfer propositions.
\end{document}
